\section{webwork-check and webwork-deploy}\label{sec:ref-webwork}

The two WeBWorK scripts moved into the machinery on 2026-09-06 and were not
documented at the time. This entry is written retrospectively, so it records
what they do \emph{now} rather than how they got here.

Both take the course repo to be the \emph{parent} of the machinery directory,
overridable with \env{COURSE\_ROOT}, and both refuse to run if that parent has
no \file{.course-machinery/} in it.

\subsection*{Where a course keeps its set definitions}

A \file{.def} file is course curation --- which problems, in which order, for
which course --- while a topic repo owns the problems themselves and is meant
to be reusable by a second course that writes its own sets. So the
\file{.def} files sit at \emph{course} level, and both scripts resolve that
directory the same way:

\begin{center}
\begin{tabularx}{\linewidth}{@{}l Y@{}}
\env{WEBWORK\_SETS\_DIR} & wins outright if set.\\
\file{WEBWORK/} & the current name, used when the directory exists.\\
\file{\_assessments/webwork-assignments/} & the pre-2026-09-10 name, still
  honoured so a course that has not moved keeps working.\\
\end{tabularx}
\end{center}

The rename followed the audience-first scheme the course roots already use
(\file{SYLLABUS/}, \file{TAs/}, \file{ADMIN/}, \file{TESTS/}), which a
lowercase \file{\_assessments/} with a nested second level fitted badly.
\emph{The fallback is temporary}: drop it once every course has moved, and the
lookup collapses back to one line. MATH1003 has moved; MATH3503 has not.

\subsection*{webwork-check}

\begin{manexsource}
webwork-check [--paths] [--drift] [TOPIC ...]
\end{manexsource}

\begin{center}
\begin{tabularx}{\linewidth}{@{}l Y@{}}
\texttt{-{}-paths} & every \texttt{source\_file} line in every \file{.def}
  points at a problem that exists. Finding \emph{no} \file{.def} files is
  itself a failure --- a check that cannot find its input must say so.\\
\texttt{-{}-drift} & every \file{.pg} under a topic's \file{webwork/diffyqs/}
  and \file{webwork/opl/} either matches the clone it was copied from, or
  records the blob it was copied from and is checked against that.\\
\end{tabularx}
\end{center}

With no flag, both run. A \env{TOPIC} operand narrows \texttt{-{}-drift} only:
the sets live in one course-level directory and each names its own topic, so
there is nothing per-topic for \texttt{-{}-paths} to select.

\paragraph{The OPL has two problem trees.}
This is the trap worth knowing. The Open Problem Library clone holds
\file{OpenProblemLibrary/} (the curated library, served at
\file{templates/Library/}) and \file{Contrib/} (contributed collections, served
at \file{templates/Contrib/}) as \emph{siblings}. A single clone variable
pointing at \file{OpenProblemLibrary/} therefore cannot see \file{Contrib/} at
all, which reported every contributed problem as \texttt{GONE} and every
\texttt{source\_file = Contrib/\ldots} line as \texttt{UNKNOWN SOURCE}. Fixed
2026-09-10, when MATH1003 began harvesting the CCCS collection.

\begin{center}
\begin{tabularx}{\linewidth}{@{}l Y@{}}
\env{OPL\_ROOT} & the clone root; moves the two below together.\\
\env{OPL\_CLONE} & \file{OpenProblemLibrary/} within it. Unchanged in meaning,
  so existing overrides still work.\\
\env{OPL\_CONTRIB} & \file{Contrib/} within it.\\
\env{DIFFYQS\_CLONE} & Lebl's \file{diffyqs-webwork} clone.\\
\end{tabularx}
\end{center}

\paragraph{Provenance is the second path level.}
A topic's \file{webwork/} has one subdirectory per upstream, named for it:

\begin{center}
\begin{tabularx}{\linewidth}{@{}l Y@{}}
\file{opl/} & from \file{OpenProblemLibrary/}, the curated tree.\\
\file{contrib/} & from \file{Contrib/}, its \emph{sibling} upstream --- not
  part of the OPL proper, and not reviewed.\\
\file{diffyqs/} & from Lebl's \file{diffyqs-webwork}.\\
\file{local/} & ours outright.\\
\end{tabularx}
\end{center}

Until 2026-09-11 the contributed copies sat under \file{opl/Contrib/}, which
reads as ``the Contrib part of the OPL'' --- a tree that does not exist ---
and which flattened away \file{OpenProblemLibrary/} for one upstream while
keeping \file{Contrib/} for the other. The sibling layout matches the clone,
matches \file{templates/Library} and \file{templates/Contrib} on the server,
and removes the special case drift-checking needed to undo the nesting.

Post-flip \file{.def} paths need no special case either way, being
topic-rooted already.

\paragraph{Drift outcomes.}
\texttt{DIFFERS} (differs from upstream with no blob recorded --- unexpected),
\texttt{MOVED} (upstream changed a file we deliberately modified --- news, not
failure), \texttt{GONE} (upstream deleted or renamed it), \texttt{AHEAD}
(upstream has problems in these sections that we do not, reported as a count).

An edited problem \emph{stays where it is} and records its origin in a header
line:

\begin{manexsource}
# UPSTREAM-BLOB: 02b6a97499f4a05af144116dacfc0835c9bcc0ec
\end{manexsource}

That separates the two questions drift bundles together --- ``have I edited
this?'' and ``has upstream moved?'' --- and keeps the second answerable on
exactly the files where an upstream fix is most expensive to absorb by hand.
Filing edited problems away under \file{local/} would give that up;
\file{local/} is for problems that are ours outright.

\subsection*{webwork-deploy}

\begin{manexsource}
webwork-deploy [-o DIR] [-n] [TOPIC ...]
\end{manexsource}

Builds one uploadable \file{.tgz} per topic, to be unpacked inside the
course's \file{templates/} on the server, producing a topic-rooted tree:

\begin{manexsource}
templates/<topic>/opl/Rochester/setAlgebra15Functions/x.pg
templates/<topic>/contrib/CCCS/CollegeAlgebra/3.1/x.pg
templates/<topic>/local/...
templates/<topic>/LICENSE.md
\end{manexsource}

The \file{.def} files go in at the top level of \file{templates/} separately.
The tarball root is named for the \emph{topic}, not the repo directory, which
is why it stages into a temporary directory rather than tarring in place.

\begin{center}
\begin{tabularx}{\linewidth}{@{}l Y@{}}
\texttt{-o DIR} & write tarballs here (default \file{./\_deploy}, gitignored).\\
\texttt{-n} & dry run: report what would go in, build nothing.\\
\end{tabularx}
\end{center}

\emph{Upload the tarball before importing the sets.} Post-flip paths point
inside \file{templates/<topic>/}, so importing first leaves every problem in
the set unresolvable.

\texttt{webwork-deploy} runs \texttt{webwork-check} first and refuses to build
if it fails --- shipping a set whose paths do not resolve is the failure mode
this layout exists to prevent.
